% !TEX program = pdflatex
\documentclass[11pt]{article}

% ── Packages ──────────────────────────────────────────────────────────
\usepackage[margin=1in]{geometry}
\usepackage{amsmath,amssymb,amsthm}
\usepackage{booktabs}
\usepackage{graphicx}
\usepackage{hyperref}
\usepackage{cleveref}
\usepackage{xcolor}
\usepackage{caption}
\usepackage{subcaption}
\usepackage{multirow}
\usepackage{enumitem}
\usepackage[numbers]{natbib}
\bibliographystyle{plainnat}

\hypersetup{colorlinks=true,linkcolor=blue!60!black,citecolor=green!50!black,urlcolor=blue!70!black}

% ── Macros ────────────────────────────────────────────────────────────
\newcommand{\Einf}{E_\infty}
\newcommand{\Reval}{R_{\mathrm{eval}}}
\newcommand{\Rtrain}{R_{\mathrm{train}}}
\newcommand{\relLtwo}{\varepsilon_{\mathrm{rel}}}
\DeclareMathOperator*{\argmin}{arg\,min}

% ── Auto-populated result macros (Phase 0–5 values) ──────────────────
\newcommand{\totalRuns}{32{,}783}
\newcommand{\mainRuns}{31{,}370}
\newcommand{\ablationRuns}{1{,}413}
\newcommand{\diagCells}{3{,}374}
\newcommand{\crossResCells}{3{,}360}
\newcommand{\crossResRows}{13{,}440}
\newcommand{\nDiagTraining}{17{,}930}
% Regime distribution (Phase 3a)
\newcommand{\regMixed}{1{,}641}
\newcommand{\regDataLim}{398}
\newcommand{\regResSens}{387}
\newcommand{\regCapLim}{186}
\newcommand{\regUnstable}{81}
\newcommand{\regSaturated}{51}
\newcommand{\regAmbiguous}{483}
\newcommand{\regDataBug}{147}
% Calibrated thresholds
\newcommand{\threshLow}{0.037}
\newcommand{\threshHigh}{0.227}
% Transfer (Phase 4)
\newcommand{\fnoWorstDeg}{6.56}
\newcommand{\fnoMeanDeg}{2.59}
\newcommand{\deeponetWorstDeg}{1.49}
\newcommand{\darcyFnoWorstDeg}{1.03}
% Convergence (Phase 5)
\newcommand{\sgdFnoRatio}{3.0}
% Legacy (for appendix reuse)
\newcommand{\holdoutWinFraction}{15/33 (45\%)}
\newcommand{\holdoutMeanImprovement}{1.2\%}
\newcommand{\withinSliceUnstableFrac}{22/27 (81\%)}

% ══════════════════════════════════════════════════════════════════════
\title{%
  Response-Surface Diagnostics for Operator-Learning Scaling%
}
\author{Lucas Perrier}
\date{\today \quad (Working Draft --- v4)}

\begin{document}
\maketitle

% ══════════════════════════════════════════════════════════════════════
\begin{abstract}
Empirical scaling laws have guided resource allocation in large-scale
machine learning, yet applying them to operator learning is complicated
by a third axis---discretization resolution---that interacts with
dataset size and model capacity in ways that 1D power-law fits cannot
capture. We conduct a systematic study across three 1D PDE benchmarks
(Burgers, Darcy, diffusion), four architectures (MLP, DeepONet, FNO,
and a capacity-controlled MLP variant), seven capacity levels, and
\totalRuns{} experimental configurations, sweeping dataset size $N$,
training resolution $\Rtrain$, and evaluation resolution $\Reval$
simultaneously.

Rather than selecting a ``winning'' parametric law per 1D slice, we
characterize each cell of the $(N, \Rtrain, D)$ grid by its
\emph{local sensitivity} to each axis and assign it a regime label:
data-limited, resolution-sensitive, capacity-limited, saturated, or
unstable. The resulting \emph{response-surface} view reveals that
48.6\% of cells are simultaneously responsive to multiple axes
(``mixed'')---structure that any 1D slice analysis collapses away.

We find three headline results. First, \textbf{the marginal value of
doubling resolution is non-monotone}: on Burgers, the only clearly
beneficial transition is $R{=}64 \to 128$ (21\% median error
reduction); going beyond $R{=}256$ provides no benefit, and 23.4\%
of all resolution-doubling transitions \emph{increase} error.
Second, \textbf{FNO's resolution invariance is task-dependent,
not architecture-inherent}: on Burgers, training at $R{=}256$
and evaluating at $R{=}32$ degrades error by
$\fnoWorstDeg{\times}$ (median worst case), while on Darcy the
worst transfer degradation is $\darcyFnoWorstDeg{\times}$.
Third, \textbf{every training divergence in the Burgers grid traces
to a data-generator bug at a single resolution ($R{=}48$)}, not to
optimizer instability---a finding that the response-surface framing
surfaced and that 1D law fits would have averaged away.

We provide per-cell regime maps, cross-resolution transfer matrices,
and practical design rules (e.g., ``use Adam for FNO---SGD converges
but to a $\sgdFnoRatio{\times}$ worse solution'') derived directly
from the response-surface structure.
\end{abstract}


% ══════════════════════════════════════════════════════════════════════
\section{Introduction}
\label{sec:intro}

Empirical scaling laws have become a central tool for understanding and
predicting the performance of large-scale machine learning
systems~\citep{kaplan2020scaling,hoffmann2022training}. In natural
language processing and computer vision, power-law relationships between
test loss and dataset size, parameter count, and compute budget guide
resource allocation decisions worth millions of dollars.

Scientific machine learning (SciML) faces an analogous resource
allocation problem with a twist: the practitioner must additionally
choose the \emph{discretization resolution} at which training data are
generated. Finer resolution is more expensive to produce but may
capture physics that coarser grids miss. For operator
learning---methods that learn maps between function
spaces~\citep{lu2021learning,li2021fourier}---the resolution axis is
especially consequential because architectures like the Fourier Neural
Operator (FNO) are marketed as ``resolution-invariant,'' yet the
degree to which this invariance holds in practice has not been
systematically studied.

The standard methodology for studying these trade-offs fits a
power-law (or, more recently, a parametric law family) to 1D slices
of error vs.\ one scaling axis at a time. This produces an exponent
or a ``winning law identity'' per slice, which is then interpreted
as a characteristic of the task--architecture pair. We argue this
framing has three limitations:

\begin{enumerate}
  \item \textbf{It collapses 2D/3D structure into 1D.} The decision
    of which axis to slice along biases the result; the joint $(N, R, D)$
    response surface contains information that no single slice retains.
  \item \textbf{It depends on the candidate law set.} Adding or
    removing a candidate can flip the ``winner'' without any change
    in the data.
  \item \textbf{It is statistically fragile at per-slice sample
    sizes of 4--7.} In a prior version of this work, a held-out
    prediction contest showed that multi-law fitting matches but does
    not meaningfully outperform a power-law-only baseline
    (\holdoutWinFraction{} wins, ${\sim}$\holdoutMeanImprovement{}
    mean improvement).
\end{enumerate}

\paragraph{This paper.}
We reframe the scaling-law question from ``which parametric law
wins?'' to ``where on the $(N, R, D)$ response surface is the
system operating?'' For each cell of the joint grid, we compute
local sensitivities to each axis, assign a regime label
(data-limited, resolution-sensitive, capacity-limited, saturated,
or unstable), and study how models transfer across resolutions via
transfer matrices. This produces actionable diagnostics---e.g.,
``your FNO is data-limited at this $(N, R)$ cell; add more training
data rather than increasing resolution''---that law-identity claims
do not.

\paragraph{Contributions.}
\begin{enumerate}
  \item We introduce \textbf{response-surface diagnostics} for
    operator learning: per-cell regime labels derived from local
    log-log sensitivities, with empirically calibrated thresholds.
    We show that 48.6\% of cells in our grid are ``mixed''
    (simultaneously responsive to multiple axes), quantifying the
    information loss inherent in 1D-slice analyses.
  \item We provide the first systematic \textbf{cross-resolution
    transfer analysis} for operator learning, showing that FNO's
    claimed resolution invariance holds on Darcy
    ($\darcyFnoWorstDeg{\times}$ worst degradation) but fails
    catastrophically on Burgers ($\fnoWorstDeg{\times}$), with a
    consistent train-high-eval-low asymmetry.
  \item We show that \textbf{the marginal value of resolution is
    non-monotone}: on Burgers, the only clearly beneficial
    transition is $R{=}64 \to 128$; beyond $R{=}256$, more
    resolution provides no benefit. Overall, 23.4\% of
    resolution-doubling transitions \emph{increase} error.
  \item We identify and trace a \textbf{data-generator bug at
    $R{=}48$} that accounts for all 582 ``training divergences''
    previously attributed to optimizer instability, demonstrating
    that the response-surface framing surfaces failure modes that
    1D analyses average away.
  \item We provide \textbf{practical design rules} grounded in the
    response-surface structure: optimizer choice (Adam vs.\ SGD:
    $\sgdFnoRatio{\times}$ error difference for FNO), the
    resolution sweet spot ($R \in [64, 128]$ for Burgers-type
    problems), and the resolution--capacity confound resolution
    via a controlled-parameter MLP.
\end{enumerate}


% ══════════════════════════════════════════════════════════════════════
\section{Related Work}
\label{sec:related}

\paragraph{Neural scaling laws.}
\citet{kaplan2020scaling} and \citet{hoffmann2022training} established
power-law scaling of cross-entropy loss with dataset tokens, parameter
count, and compute for autoregressive language models.
\citet{hestness2017deep} showed similar trends across vision and speech.
In scientific ML, \citet{perrier2025scaling} introduced scaling-law
diagnostics for physics-informed approaches (PINNs, HNNs) in
low-dimensional settings, finding that physics priors shift the error
floor rather than the scaling exponent. Our work extends this program to
infinite-dimensional operator learning and elevates resolution to a
first-class scaling variable.

\paragraph{Operator learning.}
DeepONet~\citep{lu2021learning} builds on the universal approximation
theorem for operators~\citep{chen1995universal}, factorizing the learned
operator into branch and trunk components. The Fourier Neural Operator
(FNO)~\citep{li2021fourier} parameterizes integral kernels via
truncated Fourier series, achieving state-of-the-art accuracy on
several PDE benchmarks. Recent extensions include U-shaped
architectures~\citep{wen2022ufno}, attention
mechanisms~\citep{hao2023gnot}, and physics-informed
variants~\citep{wang2021learning}.

\paragraph{Resolution and discretization in operator learning.}
A key selling point of neural operators is their ability to generalize
across resolutions. \citet{li2021fourier} demonstrated FNO's
zero-shot super-resolution capability. However, systematic studies of
how \emph{training} resolution affects accuracy---and whether more
resolution is always better---are lacking. Our transfer-matrix analysis
fills this gap and provides the first quantitative evidence that FNO's
resolution invariance is task-dependent.


% ══════════════════════════════════════════════════════════════════════
\section{Experimental Framework}
\label{sec:framework}

\subsection{Problem Setting}
\label{sec:problem}

We consider the supervised learning of solution operators
$\mathcal{G}^\dagger : \mathcal{U} \to \mathcal{V}$, where
$\mathcal{U}$ and $\mathcal{V}$ are Banach spaces of functions on
some domain $\Omega \subset \mathbb{R}^d$. Given $N$ input--output
pairs $\{(u_i, v_i)\}_{i=1}^N$ with
$v_i = \mathcal{G}^\dagger(u_i)$, a model $\mathcal{G}_\theta$ with
$D = |\theta|$ parameters is trained on data discretized at spatial
resolution $R$ (grid points per spatial dimension).

\subsection{PDE Benchmarks}

\begin{table}[h]
\centering
\caption{PDE benchmark tasks. All tasks map an input function to an
output function on the same domain, discretized at resolution $R$.}
\label{tab:tasks}
\begin{tabular}{@{}llll@{}}
\toprule
\textbf{Task} & \textbf{PDE} & \textbf{Operator} & \textbf{Character} \\
\midrule
Burgers & $u_t + u u_x = \nu u_{xx}$ &
  $u_0(x) \mapsto u(x, T)$ & Advection-dominated \\
Darcy & $-\nabla \cdot (a \nabla u) = f$ &
  $a(x) \mapsto u(x)$ & Elliptic, multiscale \\
Diffusion & $u_t = \nu u_{xx}$ &
  $u_0(x) \mapsto u(x, T)$ & Diffusion-dominated \\
\bottomrule
\end{tabular}
\end{table}

\noindent
Burgers uses a spectral solver with integrating-factor RK4 and
$\frac{2}{3}$-rule dealiasing ($\Delta t = 10^{-3}$).
Darcy uses a spectral Galerkin method.
Diffusion uses an exact spectral solution. All solvers are
resolution-parameterized.

\subsection{Architectures}

\begin{table}[h]
\centering
\caption{Model architectures and their resolution handling.}
\label{tab:models}
\begin{tabular}{@{}llp{7cm}@{}}
\toprule
\textbf{Model} & \textbf{Type} & \textbf{Resolution handling} \\
\midrule
MLP & Discretization-tied &
  Flattened input of dimension $R$; parameter count $\propto R$ \\
MLP-ctrl & Discretization-tied &
  Hidden width adjusted so $D \approx \text{const}$ across $R$;
  isolates resolution from capacity \\
DeepONet & Mesh-free &
  Branch encodes input at training $R$; trunk evaluates at arbitrary
  query points \\
FNO & Spectral &
  Global convolutions in Fourier space; truncation to $k_{\max}$ modes
  enables resolution transfer \\
\bottomrule
\end{tabular}
\end{table}

\noindent
MLP-ctrl (\texttt{mlp\_controlled}) exists on Burgers and Diffusion
only. It uses the same architecture as the standard MLP but adjusts the
hidden width at each resolution so that the parameter count matches the
standard MLP at $R{=}64$. This decouples the resolution and capacity
axes, which are confounded in the standard MLP (where $D \propto R$).
Regime-label analysis confirms the decoupling works: MLP-ctrl has only
15 resolution-sensitive cells out of 112 on Burgers, versus 107/602 for
the standard MLP (\Cref{sec:regimes}).

\subsection{Experimental Grid}

We sweep over five axes: dataset size $N$ (11 values for Burgers/Darcy,
5 for Diffusion), training resolution $\Rtrain$ (10 values for Burgers,
4 for Darcy/Diffusion), model capacity (7 named levels from
\emph{tiny} to \emph{xlarge}), data seed (2), and training seed
(2--3 per task). The total is \totalRuns{} training and evaluation
records (\mainRuns{} main grid + \ablationRuns{} optimizer ablation).

\begin{table}[h]
\centering
\caption{Experimental run accounting.}
\label{tab:run-accounting}
\small
\begin{tabular}{@{}llrl@{}}
\toprule
\textbf{Component} & \textbf{Scope} & \textbf{Runs} &
\textbf{Purpose} \\
\midrule
Main grid (diagonal) & 3 tasks $\times$ 4 models & \nDiagTraining{} &
  Training runs across $(N, R, C)$ \\
Cross-resolution eval & DeepONet + FNO & \crossResRows{} &
  Post-hoc $\Reval \neq \Rtrain$ evaluation \\
Adam $2\times$ patience & Burgers $\times$ 3 models & 756 &
  Optimizer ablation \\
SGD + cosine & Burgers $\times$ 3 models & 657 &
  Optimizer ablation \\
\midrule
\textbf{Grand total} & & \textbf{\totalRuns{}} & \\
\bottomrule
\end{tabular}
\end{table}

\paragraph{Known data-generator artifact.}
The Burgers spectral solver produces NaN/Inf in 1--4\% of generated
samples when $R{=}48$, causing all 582 training runs at that resolution
to fail before completing the first epoch. Every other resolution
(including $R{=}16, 32, 64, 96, 128, 192, 256, 384, 512$) has zero
divergences across all models, capacities, and dataset sizes. We
exclude $R{=}48$ from all primary analyses and report it separately
(\Cref{app:r48}).

\subsection{Cross-Resolution Evaluation}

After training, each DeepONet and FNO model is evaluated at all
resolutions $\Reval \in \{32, 64, 128, 256, 512\}$ except its
training resolution, yielding \crossResRows{} additional evaluation
records. MLP models are discretization-tied and cannot be evaluated
at $\Reval \neq \Rtrain$. For DeepONet, branch inputs at non-training
resolutions are linearly interpolated to $\Rtrain$ before the forward
pass. FNO natively handles arbitrary input resolutions via Fourier
truncation/zero-padding.

\subsection{Primary Outcomes and Regime Labels}
\label{sec:outcomes}

For each diagonal cell $(T, M, C, N, \Rtrain)$, we aggregate over
seeds and compute: (1)~\emph{median} test relative $L^2$ error
$\relLtwo$ (robust to outlier seeds); (2)~interquartile range;
(3)~seed coefficient of variation; (4)~convergence rate (fraction
of seeds that completed training); (5)~divergence and instability
rates.

\paragraph{Local sensitivities.}
At each cell, we compute centered log--log finite differences against
each axis:
\begin{equation}
  S_X(i) = \frac{\ln E(X_{i+1}) - \ln E(X_{i-1})}
                 {\ln X_{i+1} - \ln X_{i-1}}, \quad
  X \in \{N, R, D\}.
\end{equation}
Boundary cells use one-sided differences. The capacity axis uses
parameter count $D$ (not the named slot) as the coordinate.

\paragraph{Regime labels.}
We pool $|S_N|$, $|S_R|$, $|S_D|$ across all non-$R{=}48$ cells
and set the thresholds at the empirical quartiles:
$\theta_{\mathrm{low}} = \threshLow$ (25th percentile) and
$\theta_{\mathrm{high}} = \threshHigh$ (75th percentile). A cell is:
\begin{itemize}[nosep]
  \item \textbf{data-limited} if $|S_N| > \theta_{\mathrm{high}}$
    and $|S_R|, |S_D| < \theta_{\mathrm{high}}$;
  \item \textbf{resolution-sensitive} if $|S_R| > \theta_{\mathrm{high}}$
    and $|S_N|, |S_D| < \theta_{\mathrm{high}}$;
  \item \textbf{capacity-limited} if $|S_D| > \theta_{\mathrm{high}}$
    and $|S_N|, |S_R| < \theta_{\mathrm{high}}$;
  \item \textbf{saturated} if all three $|S| < \theta_{\mathrm{low}}$;
  \item \textbf{unstable} if divergence rate $> 0.10$ or seed CV
    $> 0.50$ (takes precedence);
  \item \textbf{mixed} otherwise (multiple axes simultaneously
    responsive).
\end{itemize}
The thresholds are not hand-picked; they fall on quartile boundaries
of the empirical sensitivity distribution and are fully reproducible
from the data.


% ══════════════════════════════════════════════════════════════════════
\section{Results: Response-Surface Diagnostics}
\label{sec:regimes}

\subsection{Regime Distribution}

\Cref{tab:regime-dist} shows the distribution of regime labels across
all \diagCells{} diagonal cells.

\begin{table}[h]
\centering
\caption{Regime label distribution across the $(N, R, D)$ grid.
``Mixed'' means $\geq$2 axes have sensitivity above $\theta_{\mathrm{high}}$.}
\label{tab:regime-dist}
\begin{tabular}{@{}lrr@{}}
\toprule
\textbf{Regime} & \textbf{Count} & \textbf{Share} \\
\midrule
mixed                & \regMixed{}   & 48.6\% \\
data-limited         & \regDataLim{} & 11.8\% \\
resolution-sensitive & \regResSens{} & 11.5\% \\
capacity-limited     & \regCapLim{}  & 5.5\% \\
unstable             & \regUnstable{} & 2.4\% \\
saturated            & \regSaturated{} & 1.5\% \\
\midrule
ambiguous (boundary) & \regAmbiguous{} & 14.3\% \\
data-bug ($R{=}48$)  & \regDataBug{} & 4.4\% \\
\bottomrule
\end{tabular}
\end{table}

\noindent
\textbf{The 48.6\% ``mixed'' fraction is the headline finding.}
Almost half of the cells in the operator-learning grid are
simultaneously responsive to more than one of $(N, R, D)$.
Forcing each cell to commit to a single ``winning law'' along a
single axis---as the standard 1D-slice methodology does---discards
the joint information in these cells.

\subsection{Per-Architecture Regime Patterns}
\label{sec:per-arch}

The regime distribution varies by (task, model), providing
architecture-specific diagnostics (\Cref{tab:regime-by-model}).

\begin{table}[h]
\centering
\caption{Per-(task, model) regime counts (excluding
ambiguous and $R{=}48$). Columns: CL = capacity-limited,
DL = data-limited, M = mixed, RS = resolution-sensitive,
S = saturated, U = unstable.}
\label{tab:regime-by-model}
\small
\begin{tabular}{@{}llrrrrrr@{}}
\toprule
\textbf{Task} & \textbf{Model} & CL & DL & M & RS & S & U \\
\midrule
Burgers & DeepONet      & 36 &  3 & 233 & 117 & 3 & 0 \\
Burgers & FNO           &  3 & 91 & 206 &  92 & 0 & 0 \\
Burgers & MLP           & 14 & 35 & 235 & 107 & 1 & 0 \\
Burgers & MLP-ctrl      &  5 &  1 &  91 &  15 & 0 & 0 \\
Darcy   & DeepONet      & 34 & 52 & 157 &  15 & 6 & 44 \\
Darcy   & FNO           &  4 & 28 & 247 &   3 & 26 & 0 \\
Darcy   & MLP           &  9 & 77 & 173 &   2 & 11 & 36 \\
Diff.   & DeepONet      & 24 & 24 &  89 &   2 & 1 & 0 \\
Diff.   & FNO           & 28 & 44 &  61 &   3 & 3 & 1 \\
Diff.   & MLP           & 17 & 41 &  69 &  13 & 0 & 0 \\
Diff.   & MLP-ctrl      & 12 &  2 &  80 &  18 & 0 & 0 \\
\bottomrule
\end{tabular}
\end{table}

Key observations:
\begin{itemize}
  \item \textbf{FNO on Burgers is the most data-limited model}
    (91 data-limited cells): FNO needs more data, not higher
    resolution, at most operating points.
  \item \textbf{DeepONet on Burgers has the most resolution-sensitive
    cells} (117): DeepONet is the most brittle to the choice of
    training resolution.
  \item \textbf{MLP-ctrl decouples resolution from capacity.}
    Only 15 of 112 Burgers/MLP-ctrl cells are resolution-sensitive
    (13.4\%), versus 107/602 (17.8\%) for the standard MLP.
  \item \textbf{Darcy instability is model-specific.} Darcy/DeepONet
    (44 unstable) and Darcy/MLP (36 unstable) account for 80 of 81
    main-grid unstable cells. Darcy/FNO has zero. This is a task--model
    interaction, not a generic instability.
  \item \textbf{Darcy/FNO has the most saturated cells} (26): FNO
    finds Darcy easy and hits a plateau where more data, resolution,
    or capacity does not help.
\end{itemize}


\subsection{Marginal Value of Resolution}
\label{sec:marginal-R}

For each cell where both $R$ and $2R$ exist with successful runs,
we compute the ratio $\relLtwo(2R) / \relLtwo(R)$. Ratio $< 1$
means doubling resolution helps; ratio $> 1$ means it hurts.

\begin{table}[h]
\centering
\caption{Median ratio of error after doubling $R$, by (task, model).
Values $>1$ indicate doubling $R$ hurts. Bold: most notable entries.}
\label{tab:marginal-R}
\begin{tabular}{@{}llrr@{}}
\toprule
\textbf{Task} & \textbf{Model} & \textbf{$n$ pairs} &
\textbf{Med.\ ratio} \\
\midrule
Burgers & FNO        & 427 & \textbf{0.878} \\
Burgers & DeepONet   & 427 & 0.956 \\
Burgers & MLP        & 427 & 0.953 \\
Burgers & MLP-ctrl   &  84 & 0.955 \\
Darcy   & FNO        & 231 & 0.995 \\
Darcy   & DeepONet   & 231 & 0.989 \\
Darcy   & MLP        & 231 & 0.974 \\
Diff.   & MLP-ctrl   &  84 & \textbf{1.213} \\
Diff.   & MLP        & 105 & \textbf{1.095} \\
Diff.   & DeepONet   & 105 & 1.041 \\
Diff.   & FNO        & 105 & 0.998 \\
\bottomrule
\end{tabular}
\end{table}

\noindent
Across the entire grid, \textbf{23.4\% of resolution-doubling
transitions increase error by $>$5\%}. For Diffusion/MLP-ctrl,
the median ratio is 1.21---doubling $R$ makes error 21\% worse.
For Burgers/FNO, doubling $R$ helps (median 0.878), but
\Cref{tab:burgers-transitions} shows the benefit is concentrated
in the $R{=}64 \to 128$ transition; beyond $R{=}256$, the marginal
value is zero.

\begin{table}[h]
\centering
\caption{Median error ratio for each resolution-doubling transition
on Burgers (all models combined).}
\label{tab:burgers-transitions}
\begin{tabular}{@{}ccr@{}}
\toprule
$R_{\text{low}}$ & $R_{\text{high}}$ & Median ratio \\
\midrule
16  & 32  & 1.855 \\
32  & 64  & 0.924 \\
64  & 128 & \textbf{0.788} \\
96  & 192 & 0.868 \\
128 & 256 & 0.957 \\
192 & 384 & 1.001 \\
256 & 512 & 1.008 \\
\bottomrule
\end{tabular}
\end{table}

\noindent
The $R{=}16 \to 32$ entry (1.855) means the typical cell gets nearly
twice as bad when going from $R{=}16$ to $R{=}32$. This reflects
that $R{=}16$ is so coarse that the test problem is artificially
easy (the ground truth itself is heavily smoothed). Once the grid
resolves the actual dynamics, the model must work harder.

\paragraph{Practical design rule.}
For Burgers-type advection-dominated problems, the sweet spot is
$R \in [64, 128]$. Below $R{=}32$, error metrics are misleading;
above $R{=}256$, additional resolution is wasted compute.


% ══════════════════════════════════════════════════════════════════════
\section{Results: Cross-Resolution Transfer}
\label{sec:transfer}

Only DeepONet and FNO support cross-resolution evaluation (MLP models
are discretization-tied). For each $(T, M, C, N)$ slice, we construct
a transfer matrix $K[i,j] = \text{median } \relLtwo$ at
$(\Rtrain = R_i, \Reval = R_j)$ and compute the degradation ratio
$K[i,j] / K[i,i]$ relative to the diagonal baseline
(\Cref{tab:transfer-scores}).

\begin{table}[h]
\centering
\caption{Per-(task, model) median transfer scores across all slices.
Higher degradation = more fragile.}
\label{tab:transfer-scores}
\begin{tabular}{@{}llrrr@{}}
\toprule
\textbf{Task} & \textbf{Model} & \textbf{Mean deg.} &
\textbf{Worst deg.} & \textbf{Asymmetry} \\
\midrule
Burgers & DeepONet & 1.015 & \deeponetWorstDeg{} & 1.44 \\
Burgers & FNO      & \fnoMeanDeg{} & \textbf{\fnoWorstDeg{}} & 1.65 \\
Darcy   & DeepONet & 1.039 & 1.13 & 0.98 \\
Darcy   & FNO      & 1.002 & \textbf{\darcyFnoWorstDeg{}} & 1.00 \\
Diff.   & DeepONet & 1.072 & 1.22 & 0.98 \\
Diff.   & FNO      & 1.074 & 1.32 & 1.02 \\
\bottomrule
\end{tabular}
\end{table}

\subsection{FNO is Catastrophically Fragile on Burgers}

All 10 worst transfer slices in the dataset are Burgers/FNO, training
at $R \in \{128, 256\}$ and evaluating at $R{=}32$, with 8--9$\times$
degradation. The architecture marketed as ``resolution-invariant''
shows the worst resolution fragility in our study. The degradation
scales super-linearly with the $\Rtrain / \Reval$ ratio.

\subsection{Transfer is Task-Dependent, Not Architecture-Inherent}

The same FNO that degrades by $\fnoWorstDeg{\times}$ on Burgers
degrades by only $\darcyFnoWorstDeg{\times}$ on Darcy. The transfer
property is not an architectural feature---it depends on the physics.
On the easier, smoother task (Darcy), FNO transfers nearly perfectly;
on the harder, advection-dominated task (Burgers), high-resolution
features have no spectral support at low resolution.

\subsection{Transfer Asymmetry}

On Burgers, the train-high-eval-low direction is consistently worse
than the reverse. FNO asymmetry is 1.65 (65\% worse in the
high-to-low direction); DeepONet is 1.44. On Darcy and Diffusion,
transfer is nearly symmetric ($\sim$1.0). The asymmetry is consistent
with a spectral-aliasing mechanism: features learned at high resolution
have no spectral support at low resolution, while features learned at
low resolution are merely undersampled.

\paragraph{Practical design rule.}
If deploying a Burgers-type FNO trained at $R \geq 128$, do not
evaluate at $R{=}32$---expect 8$\times$ error blow-up. Retrain at
the target resolution, or accept $\sim$1.5$\times$ degradation by
evaluating at $R{=}64$ instead.


% ══════════════════════════════════════════════════════════════════════
\section{Results: Convergence and Instability}
\label{sec:convergence}

\subsection{Main-Grid Convergence}

\textbf{Zero training runs experienced optimization divergence on
clean data.} Across all \nDiagTraining{} diagonal training runs
(excluding $R{=}48$), every model, capacity, and dataset-size
combination converged at 100\%. The 582 runs previously reported as
``diverged'' all occurred at $R{=}48$ due to a data-generator bug:
the Burgers spectral solver produces NaN in 1--4\% of samples at
that resolution only. See \Cref{app:r48} for the standalone
reproduction.

\subsection{Seed Variability}

The \regUnstable{} cells flagged as ``unstable'' (seed CV $> 0.50$)
are almost entirely Darcy: 44 Darcy/DeepONet and 36 Darcy/MLP cells.
Darcy/FNO has zero unstable cells. This is a task--model interaction:
Darcy's log-normal permeability fields produce challenging training
instances that DeepONet and MLP struggle with when the random seed is
unfavorable, but that FNO handles robustly.

\subsection{Optimizer Ablation}

We retrained all Burgers/medium-capacity cells with two alternative
optimizer configurations: SGD with cosine annealing (657 runs) and
Adam with doubled early-stopping patience (756 runs). All three
configurations---the baseline Adam, SGD-cosine, and
Adam-$2\times$---achieve 100\% convergence.

However, error magnitude differs substantially:
\begin{itemize}[nosep]
  \item SGD-cosine produces $\sgdFnoRatio{\times}$ worse error for
    FNO, 1.4$\times$ worse for DeepONet and MLP.
  \item Adam-$2\times$ matches the baseline within 4\%.
\end{itemize}
FNO's spectral convolution layers require per-parameter learning-rate
adaptation to find a good solution; the fixed learning-rate schedule
of SGD fails to explore the loss landscape effectively.

\paragraph{Practical design rule.}
If using FNO, use Adam (or similar adaptive optimizer). SGD converges
but to a $\sgdFnoRatio{\times}$ worse solution.


% ══════════════════════════════════════════════════════════════════════
\section{Discussion}
\label{sec:discussion}

\subsection{The Case for Response Surfaces}

The traditional 1D-slice approach to scaling analysis produced, in an
earlier version of this work, AICc-selected ``winning'' laws for 27
task--model--axis combinations. A held-out prediction contest showed
that the multi-law framework outperformed a power-law-only baseline
in only \holdoutWinFraction{} cases, with ${\sim}$\holdoutMeanImprovement{}
mean improvement. Within-slice bootstrap analysis found
\withinSliceUnstableFrac{} of slices had fragile law identity.

The response-surface framing replaces this with a richer, more stable
description:
\begin{itemize}
  \item 48.6\% of cells are ``mixed'' (simultaneously responsive to
    $\geq$2 axes)---information no 1D slice retains.
  \item Transfer matrices reveal 8--9$\times$ degradation that is
    invisible in diagonal-only law fits.
  \item The $R{=}48$ data bug was surfaced by the per-cell heatmap;
    it would have been averaged into a ``3.3\% divergence rate''
    under a 1D analysis.
\end{itemize}

\subsection{When Does Resolution Matter?}

Resolution matters most for advection-dominated problems (Burgers)
where it allows the model to resolve sharp gradients. It matters least
for smooth problems (Diffusion), where higher resolution introduces
irrelevant high-frequency content that MLPs waste capacity fitting.
For elliptic problems (Darcy), the effect is intermediate and
model-dependent.

The controlled-parameter MLP experiment confirms that ``resolution
hurts'' on Diffusion is a genuine resolution effect, not a
parameter-count confound: degradation persists even at fixed $D$.

\subsection{Architecture Recommendations}

\begin{enumerate}
  \item \textbf{FNO:} Most accurate overall (best cells are all
    FNO on Diffusion at $\relLtwo \approx 0.002$). But needs the most
    data (91 data-limited Burgers cells) and is the most fragile to
    cross-resolution transfer (8--9$\times$ on Burgers). Use Adam;
    SGD yields $\sgdFnoRatio{\times}$ worse solutions.
  \item \textbf{DeepONet:} More robust to transfer ($\deeponetWorstDeg{\times}$
    worst degradation on Burgers) but the most sensitive to training
    resolution (117 resolution-sensitive Burgers cells).
  \item \textbf{MLP:} Cannot evaluate at $\Reval \neq \Rtrain$
    (discretization-tied). Competitive on Darcy but worst on Burgers.
  \item \textbf{MLP-ctrl:} Successfully decouples resolution from
    capacity. Most cells are ``mixed'' rather than
    resolution-sensitive, confirming the decoupling.
\end{enumerate}

\subsection{Limitations}

\begin{enumerate}
  \item All tasks are 1D. In higher dimensions, resolution effects are
    amplified ($R^d$ cost growth) and may shift the balance.
  \item Darcy's resolution range ($R \in \{32, 64, 128, 256\}$) is
    half of Burgers's. Cross-resolution conclusions about Darcy are
    limited to this range.
  \item Diffusion uses only 5 $N$ values per model (vs.\ 11 for
    Burgers). Local-slope claims on the data axis for Diffusion are
    correspondingly weaker.
  \item We use 2 data seeds globally (3 for MLP-ctrl). Seed-variability
    estimates have limited statistical power.
  \item The regime-label thresholds are derived from the empirical
    quartiles of the sensitivity distribution. They are not universal
    constants; a different grid would yield different thresholds.
\end{enumerate}


% ══════════════════════════════════════════════════════════════════════
\section{Conclusion}
\label{sec:conclusion}

We have presented a response-surface diagnostic framework for
operator-learning scaling that replaces 1D parametric law fits with
per-cell regime labels, cross-resolution transfer matrices, and
marginal value-of-resolution analysis. The framework reveals
structure---48.6\% mixed-regime cells, task-dependent FNO fragility,
and a data-generator bug masquerading as optimizer instability---that
the standard 1D-slice methodology cannot surface.

Our findings distill into three practical design rules for
practitioners:
\begin{enumerate}
  \item \textbf{The resolution sweet spot is narrow.} For
    advection-dominated problems, $R \in [64, 128]$ is the only
    range with reliable benefit; below $R{=}32$, error metrics are
    misleading; above $R{=}256$, additional resolution is wasted
    compute.
  \item \textbf{Do not trust cross-resolution transfer on hard
    problems.} FNO trained at $R{=}256$ and evaluated at $R{=}32$
    loses 8--9$\times$ accuracy on Burgers, even though the same
    FNO transfers nearly losslessly on Darcy.
  \item \textbf{Optimizer choice matters for FNO.} SGD converges
    but to a $\sgdFnoRatio{\times}$ worse solution than Adam.
\end{enumerate}

The response-surface framework is general: it can be applied to any
experimental grid with $\geq$2 scaling axes, and the regime labels
are derived from the data without assuming any parametric form. We
release the full dataset (\totalRuns{} runs) and analysis pipeline
to support future work.

\section*{Acknowledgments}
Experiments were conducted on three RunPod GPU instances (NVIDIA A4000,
L4, and RTX 3090). Total experiment wall time was approximately 8 hours
across all pods.


% ══════════════════════════════════════════════════════════════════════
\begin{thebibliography}{20}

\bibitem[Chen and Chen(1995)]{chen1995universal}
T.~Chen and H.~Chen.
\newblock Universal approximation to nonlinear operators by neural networks
  with arbitrary activation functions and its application to dynamical systems.
\newblock \emph{IEEE Trans. on Neural Networks}, 6(4):911--917, 1995.

\bibitem[Hao et~al.(2023)]{hao2023gnot}
Z.~Hao, Z.~Wang, H.~Su, C.~Ying, Y.~Dong, S.~Liu, Z.~Cheng, J.~Song, and
  J.~Zhu.
\newblock {GNOT}: A general neural operator transformer for operator learning.
\newblock In \emph{ICML}, 2023.

\bibitem[Hestness et~al.(2017)]{hestness2017deep}
J.~Hestness, S.~Narang, N.~Ardalani, G.~Diamos, H.~Jun, H.~Kianinejad,
  M.~M.~A. Patwary, Y.~Yang, and Y.~Zhou.
\newblock Deep learning scaling is predictable, empirically.
\newblock \emph{arXiv:1712.00409}, 2017.

\bibitem[Hoffmann et~al.(2022)]{hoffmann2022training}
J.~Hoffmann, S.~Borgeaud, A.~Mensch, E.~Buchatskaya, T.~Cai, E.~Rutherford,
  D.~de~Las~Casas, L.~A. Hendricks, J.~Welbl, A.~Clark, T.~Hennigan,
  E.~Noland, K.~Millican, G.~van~den~Driessche, B.~Damoc, A.~Guy,
  S.~Osindero, K.~Simonyan, E.~Elsen, J.~W. Rae, O.~Vinyals, and
  L.~Sifre.
\newblock Training compute-optimal large language models.
\newblock In \emph{NeurIPS}, 2022.

\bibitem[Kaplan et~al.(2020)]{kaplan2020scaling}
J.~Kaplan, S.~McCandlish, T.~Henighan, T.~P. Brown, B.~Chess, R.~Child,
  S.~Gray, A.~Radford, J.~Wu, and D.~Amodei.
\newblock Scaling laws for neural language models.
\newblock \emph{arXiv:2001.08361}, 2020.

\bibitem[Li et~al.(2021)]{li2021fourier}
Z.~Li, N.~Kovachki, K.~Azizzadenesheli, B.~Liu, K.~Bhatt, A.~Stuart, and
  A.~Anandkumar.
\newblock Fourier neural operator for parametric partial differential equations.
\newblock In \emph{ICLR}, 2021.

\bibitem[Lu et~al.(2021)]{lu2021learning}
L.~Lu, P.~Jin, G.~Pang, Z.~Zhang, and G.~E. Karniadakis.
\newblock Learning nonlinear operators via {DeepONet} based on the universal
  approximation theorem of operators.
\newblock \emph{Nature Machine Intelligence}, 3(3):218--229, 2021.

\bibitem[Perrier(2025)]{perrier2025scaling}
L.~Perrier.
\newblock Scaling-law diagnostics for physics-informed machine learning.
\newblock Working paper, 2025.

\bibitem[Wang et~al.(2021)]{wang2021learning}
S.~Wang, H.~Wang, and P.~Perdikaris.
\newblock Learning the solution operator of parametric partial differential
  equations with physics-informed {DeepONets}.
\newblock \emph{Science Advances}, 7(40), 2021.

\bibitem[Wen et~al.(2022)]{wen2022ufno}
G.~Wen, Z.~Li, K.~Azizzadenesheli, A.~Anandkumar, and S.~M. Benson.
\newblock {U-FNO}---an enhanced {Fourier} neural operator-based deep-learning
  model for multiphase flow.
\newblock \emph{Advances in Water Resources}, 163:104180, 2022.

\end{thebibliography}


% ══════════════════════════════════════════════════════════════════════
\appendix

\section{Experimental Details}
\label{app:details}

\subsection{Capacity Grid}

\begin{table}[h]
\centering
\caption{MLP and DeepONet capacity grid (hidden layers $\times$ hidden
width). FNO uses a separate grid over modes, width, and number of
layers.}
\label{tab:capacity-grid}
\begin{tabular}{@{}lrrr@{}}
\toprule
\textbf{Level} & \textbf{MLP layers} & \textbf{Hidden width} &
\textbf{$\sim D$ (MLP, $R\!=\!128$)} \\
\midrule
tiny      & 2 & 32  & $\sim$5K \\
small     & 2 & 64  & $\sim$13K \\
small-med & 2 & 96  & $\sim$22K \\
medium    & 2 & 128 & $\sim$34K \\
med-large & 2 & 192 & $\sim$62K \\
large     & 2 & 256 & $\sim$82K \\
xlarge    & 3 & 256 & $\sim$148K \\
\bottomrule
\end{tabular}
\end{table}

\subsection{Training Configuration}
All models are trained with Adam (learning rate $10^{-3}$, weight decay
$10^{-5}$), batch size 64, and maximum 5{,}000 epochs with early
stopping (patience 200, monitoring validation relative $L^2$). Data is
split 80/10/10 into train/validation/test. Relative $L^2$ loss
$\|\hat{v} - v\|_2 / \|v\|_2$ is used for training.

\subsection{PDE Solver Details}

\paragraph{Burgers equation.}
$u_t + u u_x = \nu u_{xx}$ on $[0, 2\pi)$ with periodic boundary
conditions, $\nu = 0.01$, $T = 1.0$. Solved with integrating-factor
RK4 in Fourier space with $\frac{2}{3}$-rule dealiasing and
$\Delta t = 10^{-3}$. Random initial conditions are superpositions of
Fourier modes with random phases and amplitudes drawn from a $k^{-2}$
energy spectrum.

\paragraph{Darcy flow.}
$-\nabla \cdot (a(x) \nabla u(x)) = f(x)$ on $[0, 1]$ with Dirichlet
boundary conditions, $f(x) = 1$. The coefficient field $a(x)$ is drawn
from a log-normal random field. Solved via spectral Galerkin method.

\paragraph{Diffusion equation.}
$u_t = \nu u_{xx}$ on $[0, 2\pi)$ with periodic boundary conditions,
$\nu = 0.01$, $T = 1.0$. Exact spectral solution: each Fourier mode
$\hat{u}_k(t) = \hat{u}_k(0) e^{-\nu k^2 t}$. Same initial condition
distribution as Burgers.


\section{The $R{=}48$ Data-Generator Bug}
\label{app:r48}

The Burgers spectral solver (\texttt{\_spectral\_solve} in
\texttt{src/scaling\_operator\_learning/tasks/burgers.py}) produces
NaN/Inf in approximately 1--4\% of generated samples when $R{=}48$,
the only non-power-of-two resolution in the Burgers grid. All 582
training runs at $R{=}48$ fail before completing the first epoch
(\texttt{best\_epoch=-1}, \texttt{failure\_reason=nan\_or\_inf}).
Every other resolution has zero failures.

\begin{table}[h]
\centering
\caption{Standalone reproduction of the $R{=}48$ data-generator bug
(200 samples per (R, seed), no model training).}
\label{tab:r48-repro}
\begin{tabular}{@{}rrrl@{}}
\toprule
$R$ & NaN (seed=11) & NaN (seed=22) & $\max|u_T|$ \\
\midrule
16  & 0/200 & 0/200 & 4.27 \\
32  & 0/200 & 0/200 & 4.26 \\
\textbf{48} & \textbf{7/200} & \textbf{2/200} & \textbf{18.6} \\
64  & 0/200 & 0/200 & 3.14 \\
96  & 0/200 & 0/200 & 3.24 \\
128 & 0/200 & 0/200 & 3.10 \\
192 & 0/200 & 0/200 & 3.10 \\
256 & 0/200 & 0/200 & 2.94 \\
\bottomrule
\end{tabular}
\end{table}

\noindent
The 582 ``training divergences'' reported in earlier drafts are
reattributed from optimization instability to data-generator failure.
The reproducer script is at \texttt{scripts/repro\_r48\_data\_bug.py}.


\section{Law Fits as Supplementary Analysis}
\label{app:law-fits}

As a supplementary analysis, we fit six candidate functional forms
(power, exponential, logarithmic, stretched exponential, saturation,
linear) to 1D slices of error vs.\ each scaling axis, using AICc for
model selection. Full results are in the project repository at
\texttt{results/multilaw\_fits.json}.

\paragraph{Why this is not the primary analysis.}
A held-out prediction contest---the most honest test of whether
multi-law fitting adds value---found that it outperforms a
power-law-only baseline in only \holdoutWinFraction{} of cases, with
${\sim}$\holdoutMeanImprovement{} mean improvement. Within-slice
bootstrap analysis found \withinSliceUnstableFrac{} of law-identity
assignments were ``exploratory'' (bootstrap-unstable). These results
are the strongest internal evidence that law identity is not the right
level of abstraction for this dataset, and the primary motivation for
the response-surface framing adopted in this paper.

\paragraph{What the law fits do show.}
The directional findings from the law fits are consistent with the
response-surface analysis:
\begin{itemize}
  \item Data scaling consistently selects bounded-decay families
    (exponential, saturation, power with floor) over unbounded
    forms (logarithmic), once the grid is densified to 11 $N$ values.
  \item Resolution-axis law identity (logarithmic vs.\ linear) is
    robust to the candidate set and to the optimizer (\Cref{app:details}).
  \item Capacity-axis sensitivity is low across the board, consistent
    with the low $|S_D|$ values in \Cref{sec:outcomes}.
\end{itemize}


\end{document}
